\documentclass[12pt,a4paper]{report}
\usepackage{lmodern}
\usepackage[T5]{fontenc}
\usepackage[utf8]{inputenc}
\DeclareTextSymbolDefault{\DH}{T1}
\usepackage[vietnamese]{babel}
\usepackage{amsmath}
\usepackage{amssymb}
\usepackage{float}
\usepackage[left=2.5cm,right=2cm,top=2cm,bottom=2.5cm]{geometry} 
\setlength{\parindent}{0pt}
\usepackage{graphicx} 
\graphicspath{ {figures/} {Images/} }
\usepackage[table,x11names]{xcolor}
\usepackage{tikz}
\usepackage{scrextend}
\usepackage{xurl}
\usepackage[
    colorlinks=true,
    linkcolor=blue,
    urlcolor=blue,
    citecolor=blue
]{hyperref}
% \changefontsizes{14pt}
\usetikzlibrary{calc, positioning, shapes.geometric, arrows, fit, shadows}
\usepackage{verbatim}
\usepackage{enumitem}
\usepackage{array}
\usepackage{multirow}
\usepackage{listings}
\usepackage{tabularx}
\usepackage{pgfplots}
\usepackage[most]{tcolorbox}
\pgfplotsset{compat=1.18}

% Gói hỗ trợ thuật toán trong đồ án
\usepackage{algorithm,algorithmic}

% Bảng biểu bổ trợ
\newcolumntype{L}[1]{>{\raggedright\let\newline\\\arraybackslash\hspace{0pt}}m{#1}}
\newcolumntype{C}[1]{>{\centering\let\newline\\\arraybackslash\hspace{0pt}}m{#1}}
\newcolumntype{R}[1]{>{\raggedleft\let\newline\\\arraybackslash\hspace{0pt}}m{#1}}

% Tài liệu tham khảo
\usepackage[
	sorting=none,
	backend=biber,
	defernumbers=true]{biblatex}
\IfFileExists{References/references.bib}{
  \addbibresource{References/references.bib}
}{
  \IfFileExists{Project1/References/references.bib}{
    \addbibresource{Project1/References/references.bib}
  }{}
}

\makeatletter
\def\blx@maxline{77}
\makeatother

% Định nghĩa màu sắc cho code (Giống VS Code)
\definecolor{codegreen}{rgb}{0,0.6,0}
\definecolor{codegray}{rgb}{0.5,0.5,0.5}
\definecolor{codepurple}{rgb}{0.58,0,0.82}
\definecolor{backcolour}{rgb}{0.96,0.96,0.96}
\definecolor{dodo}{RGB}{153,0,0}

% Cấu hình chung (Fix lỗi tràn lề Overfull hbox)
\lstset{
    backgroundcolor=\color{backcolour},   
    commentstyle=\color{codegreen},
    keywordstyle=\color{blue},
    numberstyle=\tiny\color{codegray},
    stringstyle=\color{red},
    basicstyle=\ttfamily\footnotesize, % Dùng font nhỏ hơn để vừa trang giấy
    breakatwhitespace=false,         
    breaklines=true,                 % QUAN TRỌNG: Tự động xuống dòng
    captionpos=b,                    
    keepspaces=true,                 
    numbers=left,                    
    numbersep=5pt,                  
    showspaces=false,                
    showstringspaces=false,
    showtabs=false,                  
    tabsize=2,                       % Tab nhỏ lại để đỡ chiếm chỗ
    frame=single,                    % Khung viền đơn giản
    extendedchars=true,              % Hỗ trợ ký tự đặc biệt
    literate={á}{{\'a}}1 {à}{{\`a}}1 {ả}{{\h{a}}}1 {ã}{{\~a}}1 {ạ}{{\d{a}}}1 % Hỗ trợ tiếng Việt trong comment
}

% Định nghĩa code JavaScript
\lstdefinelanguage{JavaScript}{
  keywords={typeof, new, true, false, catch, function, return, null, catch, switch, var, if, in, while, do, else, case, break, const, let, async, await},
  keywordstyle=\color{blue}\bfseries,
  ndkeywords={class, export, boolean, throw, implements, import, this, window, document, console},
  ndkeywordstyle=\color{teal}\bfseries,
  identifierstyle=\color{black},
  sensitive=false,
  comment=[l]{//},
  morecomment=[s]{/*}{*/},
  commentstyle=\color{green!60!black}\ttfamily,
  stringstyle=\color{red}\ttfamily,
  morestring=[b]',
  morestring=[b]"
}

% Định nghĩa style cho code C#
\lstdefinelanguage{CSharp}{
  keywords={typeof, new, true, false, catch, function, return, null, catch, switch, var, if, in, while, do, else, case, break, string, byte, public, private, void, static},
  keywordstyle=\color{blue}\bfseries,
  ndkeywords={class, export, boolean, throw, implements, import, this},
  ndkeywordstyle=\color{darkgray}\bfseries,
  identifierstyle=\color{black},
  sensitive=false,
  comment=[l]{//},
  morecomment=[s]{/*}{*/},
  commentstyle=\color{green!60!black}\ttfamily,
  stringstyle=\color{red}\ttfamily,
  morestring=[b]",
  basicstyle=\small\ttfamily,
  numbers=left,
  numberstyle=\tiny\color{gray},
  stepnumber=1,
  numbersep=10pt,
  tabsize=2,
  showstringspaces=false,
  breaklines=true,
  frame=lines,
  backgroundcolor=\color{gray!5}
}

% Đổi tên mặc định
\renewcommand{\chaptername}{Chương}
\renewcommand{\figurename}{Hình}
\renewcommand{\tablename}{Bảng}
\renewcommand{\contentsname}{Mục lục}
\renewcommand{\listfigurename}{Danh sách hình}
\renewcommand{\listtablename}{Danh sách bảng}
\renewcommand{\appendixname}{Phụ lục}

% Định dạng chapter
\usepackage{titlesec}
\titleformat{\chapter}
    [display] 
    {\normalfont\bfseries\Large}{\chaptername \ \thechapter}{10pt}{\huge}
\titlespacing*{\chapter}{0pt}{-10pt}{40pt} %khoảng cách giữa chapter và đầu trang

\titleformat{\section}
    {\normalfont\bfseries\large}{\thesection}{1em}{}
 
\titleformat{\subsection}
    {\normalfont\bfseries\normalsize}{\thesubsection}{1em}{}

% --- 2. GÓI HEADER (CHỈ KHAI BÁO, CHƯA BẬT) ---
\usepackage{fancyhdr}

\begin{document}
\thispagestyle{fancy}

% --- TRANG BÌA ---
\begin{titlepage}

\begin{tikzpicture}[remember picture,overlay,inner sep=0,outer sep=0]
     \draw[blue!70!black,line width=4pt] ([xshift=-1.5cm,yshift=-2cm]current page.north east) coordinate (A)--([xshift=1.5cm,yshift=-2cm]current page.north west) coordinate(B)--([xshift=1.5cm,yshift=2cm]current page.south west) coordinate (C)--([xshift=-1.5cm,yshift=2cm]current page.south east) coordinate(D)--cycle;

     \draw ([yshift=0.5cm,xshift=-0.5cm]A)-- ([yshift=0.5cm,xshift=0.5cm]B)--
     ([yshift=-0.5cm,xshift=0.5cm]B) --([yshift=-0.5cm,xshift=-0.5cm]B)--([yshift=0.5cm,xshift=-0.5cm]C)--([yshift=0.5cm,xshift=0.5cm]C)--([yshift=-0.5cm,xshift=0.5cm]C)-- ([yshift=-0.5cm,xshift=-0.5cm]D)--([yshift=0.5cm,xshift=-0.5cm]D)--([yshift=0.5cm,xshift=0.5cm]D)--([yshift=-0.5cm,xshift=0.5cm]A)--([yshift=-0.5cm,xshift=-0.5cm]A)--([yshift=0.5cm,xshift=-0.5cm]A);

     \draw ([yshift=-0.3cm,xshift=0.3cm]A)-- ([yshift=-0.3cm,xshift=-0.3cm]B)--
     ([yshift=0.3cm,xshift=-0.3cm]B) --([yshift=0.3cm,xshift=0.3cm]B)--([yshift=-0.3cm,xshift=0.3cm]C)--([yshift=-0.3cm,xshift=-0.3cm]C)--([yshift=0.3cm,xshift=-0.3cm]C)-- ([yshift=0.3cm,xshift=0.3cm]D)--([yshift=-0.3cm,xshift=0.3cm]D)--([yshift=-0.3cm,xshift=-0.3cm]D)--([yshift=0.3cm,xshift=-0.3cm]A)--([yshift=0.3cm,xshift=0.3cm]A)--([yshift=-0.3cm,xshift=0.3cm]A);

\end{tikzpicture}

\begin{center}
    \vspace{7pt}
    
    \textbf{ĐẠI HỌC QUỐC GIA THÀNH PHỐ HỒ CHÍ MINH}
    
    \vspace{7pt}
    \textbf{TRƯỜNG ĐẠI HỌC KHOA HỌC TỰ NHIÊN}

    \vspace{7pt}
    \textbf{KHOA CÔNG NGHỆ THÔNG TIN}
\end{center}
\vspace{10pt}
\begin{center}
    \IfFileExists{hcmus.png}{\includegraphics[scale=0.2]{hcmus.png}\vspace{15pt}}{}
    
    % \fontsize{18pt}{17pt}\selectfont 
    % \textbf{Báo Cáo Đồ Án 1} 
    
    \vspace{7pt}
\end{center}
\begin{flushleft}
    \fontsize{14pt}{17pt}\selectfont  
    \textbf{\textsl{ĐỀ TÀI}}
\end{flushleft}
\begin{center}
    \fontsize{24pt}{20pt}\selectfont\textbf{{\textrm{PROJECT 2: GRIDUCTIVE SOLVER}}}
    \vspace{10pt}

    \fontsize{14pt}{17pt}\selectfont
    \textbf{CƠ SỞ TRÍ TUỆ NHÂN TẠO}
\end{center}

\begin{center}
    \vspace{12pt}
    \textbf{Lớp:} \textbf{\large 24TNT1}
    \vspace{10pt}
    
    \renewcommand{\arraystretch}{1.3}
    {\large 
    \begin{tabular}{l @{\hspace{3cm}} l} 
        \textit{\small Họ và tên} & \textit{\small MSSV} \\
        \textbf{Đinh Như Phát}      & \textbf{24122002} \\
        \textbf{Phạm Nguyên Phương} & \textbf{24122011} \\
        \textbf{Từ Văn Khôi}        & \textbf{24122019} \\
        \textbf{Nguyễn Cao Bản}     & \textbf{24122029} \\
        \textbf{Trần Lê Minh Đức}   & \textbf{24122031} \\
    \end{tabular}
    }
    \vspace{10pt}
\end{center}

\vspace{12pt}

\begin{center}
    \fontsize{14pt}{17pt}\selectfont
    \textit{Giáo viên hướng dẫn:}

    \vspace{5pt}
    
    \textbf{GS.TS. Lê Hoài Bắc} \\
    \textbf{ThS. Nguyễn Ngọc Đức}
\end{center}

\vfill

\begin{center}
    Tp. Hồ Chí Minh, ngày 22 tháng 10 năm 2026
\end{center}

\end{titlepage}


% --- MỤC LỤC ---
\pagestyle{plain} % Quan trọng: Header trắng, chỉ có số trang ở dưới
\pagenumbering{roman} % Đánh số i, ii, iii, ...
\tableofcontents
\addcontentsline{toc}{chapter}{Mục lục}

% Phần đánh giá thành viên
\IfFileExists{Appendix/danhgiatv.tex}{
  \chapter*{Đánh giá thành viên}
\phantomsection
\addcontentsline{toc}{chapter}{Đánh giá thành viên}

Bảng dưới đây tổng hợp sự phân công vai trò, khối lượng công việc phụ trách và mức độ hoàn thành nhiệm vụ của từng thành viên trong nhóm trong suốt quá trình thực hiện đồ án.

\begin{center}
\textbf{Bảng tổng hợp phân công nhiệm vụ và đánh giá mức độ hoàn thành của các thành viên}
\vspace{6pt}

\begingroup
\small
\setlength{\tabcolsep}{6pt}
\renewcommand{\arraystretch}{1.35}
\begin{tabularx}{\textwidth}{|c|l|c|X|c|}
\hline
\rowcolor{gray!15}
\textbf{STT} & \textbf{Họ và tên} & \textbf{MSSV} & \textbf{Vai trò và nội dung phụ trách chính} & \textbf{Hoàn thành} \\
\hline
1 & \textbf{Đinh Như Phát} & 24122002 & Phụ trách lõi trò chơi (\texttt{Game}), mô hình hóa trạng thái, luật di chuyển và cơ chế ô đặc biệt; điều phối tích hợp mã nguồn và xây dựng kịch bản video demo. & \textbf{100\%} \\
\hline
2 & \textbf{Phạm Nguyên Phương} & 24122011 & Hiện thực thuật toán UCS và mô hình chi phí địa hình; đối soát nghiệm tối ưu chi phí; tổng hợp bảng biểu thực nghiệm, viết phần kết luận và rà soát bản thảo báo cáo. & \textbf{100\%} \\
\hline
3 & \textbf{Từ Văn Khôi} & 24122019 & Hiện thực thuật toán A* và thiết kế hàm đánh giá heuristic; xây dựng bộ công cụ đo kiểm tự động (\texttt{benchmark runner}), vẽ biểu đồ và phân tích dữ liệu thực nghiệm. & \textbf{100\%} \\
\hline
4 & \textbf{Nguyễn Cao Bản} & 24122029 & Thiết kế và nâng cấp giao diện đồ họa Pygame, hệ thống hiệu ứng hình ảnh/âm thanh, bảng điều khiển bộ giải AI, mô phỏng bước đi và kiểm thử chấp nhận giao diện. & \textbf{100\%} \\
\hline
5 & \textbf{Trần Lê Minh Đức} & 24122031 & Hiện thực tác tử suy luận logic (\texttt{LogicAgent}), nguyên tắc không đoán mò, cơ chế gợi ý (\texttt{hint}) và truy vết suy luận; cài đặt thuật toán kiểm tra tính duy nhất (\texttt{check\_uniqueness}); xây dựng hệ thống thực nghiệm đánh giá hiệu năng (Chương 5). & \textbf{100\%} \\
\hline
\end{tabularx}
\endgroup
\end{center}

\section*{Chi tiết công việc thực hiện của từng thành viên}

\subsection*{1. Đinh Như Phát -- MSSV: 24122002}
\begin{itemize}
    \item \textbf{Nhiệm vụ chuyên môn:}
    \begin{itemize}
        \item Thiết kế cấu trúc biểu diễn trạng thái $s = (C, a, b)$ bất biến và xây dựng lõi xử lý trò chơi (\texttt{Game}), bao gồm hệ thống luật di chuyển và hàm sinh trạng thái kế tiếp.
        \item Kiểm chứng hình học các tư thế của khối hộp $1 \times 1 \times 2$ và 12 phép chuyển tiếp hướng lăn trên lưới tọa độ.
        \item Hiện thực và kiểm thử chuyên sâu các cơ chế ô nâng cao: ô dễ vỡ (\textit{fragile tiles}), cầu đóng/mở theo mặt nạ bit (\textit{bridges}), công tắc nhẹ/nặng (\textit{soft/heavy switches}) và công tắc phân tách/hợp nhất khối (\textit{split/merge}).
        \item Rà soát mã nguồn (\textit{code review}), giải quyết xung đột khi tích hợp các nhánh phát triển và chuẩn bị kịch bản, quay dựng video minh họa sản phẩm.
        \item Tham gia biên soạn Chương 1 (Giới thiệu) và Chương 3 (Hình thức hóa bài toán tìm kiếm).
    \end{itemize}
    \item \textbf{Mức độ hoàn thành:} \textbf{100\%} (Đóng góp đồng đều: 20\% khối lượng toàn đồ án).
\end{itemize}

\subsection*{2. Phạm Nguyên Phương -- MSSV: 24122011}
\begin{itemize}
    \item \textbf{Nhiệm vụ chuyên môn:}
    \begin{itemize}
        \item Nghiên cứu và hiện thực thuật toán Tìm kiếm chi phí đồng nhất (Uniform-Cost Search -- UCS) sử dụng cấu trúc hàng đợi ưu tiên Min-Heap.
        \item Thiết lập công thức chi phí địa hình $c(s, u, s') = 1 + 2F(s,s') + B(s,s')$ và kiểm chứng cơ chế cập nhật chi phí tốt nhất (\texttt{best\_cost}).
        \item Tổng hợp số liệu đo lường thực nghiệm trên 32 màn kiểm thử, thiết kế hệ thống bảng thống kê định lượng và đối chiếu độ lệch chi phí giữa các thuật toán.
        \item Biên soạn nội dung phần UCS trong Chương 4, tổng hợp bảng biểu Chương 5, viết Chương 6 (Kết luận và hướng phát triển) và kiểm tra quy chuẩn trình bày file PDF nộp cuối kỳ.
    \end{itemize}
    \item \textbf{Mức độ hoàn thành:} \textbf{100\%} (Đóng góp đồng đều: 20\% khối lượng toàn đồ án).
\end{itemize}

\subsection*{3. Từ Văn Khôi -- MSSV: 24122019}
\begin{itemize}
    \item \textbf{Nhiệm vụ chuyên môn:}
    \begin{itemize}
        \item Nghiên cứu và hiện thực thuật toán Tìm kiếm có thông tin A* (A-Star Search).
        \item Thiết kế hàm đánh giá heuristic dựa trên khoảng cách Chebyshev chia đôi, chứng minh tính chấp nhận được (\textit{admissible}) và tính nhất quán (\textit{consistent}) trên cả trạng thái nguyên khối và phân tách.
        \item Xây dựng mô-đun chạy thực nghiệm tự động (\texttt{benchmark runner}) đo thời gian thực thi bằng \texttt{time.perf\_counter} và đỉnh bộ nhớ bằng \texttt{tracemalloc}, xuất dữ liệu chuẩn hóa ra định dạng CSV/JSON.
        \item Trực quan hóa dữ liệu thành hệ thống biểu đồ so sánh 4 chỉ số hiệu năng theo từng nhóm độ khó và viết phần phân tích chuyên sâu trong Chương 4 và Chương 5.
    \end{itemize}
    \item \textbf{Mức độ hoàn thành:} \textbf{100\%} (Đóng góp đồng đều: 20\% khối lượng toàn đồ án).
\end{itemize}

\subsection*{4. Nguyễn Cao Bản -- MSSV: 24122029}
\begin{itemize}
    \item \textbf{Nhiệm vụ chuyên môn:}
    \begin{itemize}
        \item Xây dựng và hoàn thiện giao diện đồ họa người dùng (GUI) bằng thư viện Pygame, bao gồm màn hình chính, màn hình chọn màn chơi, hướng dẫn và màn hình chơi trực tiếp.
        \item Tích hợp bảng điều khiển bộ giải AI (\textit{solver panel}), chức năng phát lại lời giải từng bước (\textit{step-by-step replay}), điều chỉnh tốc độ diễn hoạt và hiển thị trực tiếp các chỉ số đo lường.
        \item Nâng cấp hệ thống trang phục khối lăn (\textit{skins}), phông nền đồng bộ, góc nhìn chiếu (\textit{camera scaling}), hiệu ứng hạt vật lý và âm thanh tổng hợp.
        \item Thực hiện kiểm thử chấp nhận thủ công (\textit{manual acceptance testing}) trên toàn bộ các màn chơi và biên soạn Chương 2 (Cài đặt trò chơi Bloxorz) cùng Phụ lục giao diện.
    \end{itemize}
    \item \textbf{Mức độ hoàn thành:} \textbf{100\%} (Đóng góp đồng đều: 20\% khối lượng toàn đồ án).
\end{itemize}

\subsection*{5. Trần Lê Minh Đức -- MSSV: 24122031}
\begin{itemize}
    \item \textbf{Nhiệm vụ chuyên môn:}
    \begin{itemize}
        \item Thiết kế và hiện thực tác tử suy luận logic tự động (\texttt{LogicAgent}) theo nguyên tắc suy luận tất định không đoán mò (\textit{no-guess deduction policy}).
        \item Xây dựng thuật toán thực thi bước suy luận đơn lẻ (\texttt{step}) với cơ chế lựa chọn ô cờ bị ép buộc tất định theo thứ tự hàng trước (\textit{row-major order}) và vòng lặp giải toàn cục (\texttt{run}).
        \item Phát triển cơ chế gợi ý nước đi thông minh (\texttt{hint}) và truy vết chuỗi dữ kiện liên quan (\texttt{clue\_ids}) thông qua hàm phân tích vùng hình học.
        \item Hiện thực thuật toán kiểm tra tính duy nhất của nghiệm ngoại tuyến (\texttt{check\_uniqueness}) độc lập không sử dụng đáp án ẩn, chứng minh tính đơn nhất chỉ với tối đa $1 + n^2$ lần gọi SAT.
        \item Thiết kế và triển khai kịch bản thực nghiệm đo lường 6 chỉ số trên toàn bộ 10 hồ sơ vụ án chuẩn và biên soạn toàn bộ Chương 5 (Thực nghiệm và Đánh giá) cùng các mục liên quan trong Chương 4.
    \end{itemize}
    \item \textbf{Mức độ hoàn thành:} \textbf{100\%} (Đóng góp đồng đều: 20\% khối lượng toàn đồ án).
\end{itemize}

}{
  \IfFileExists{Project1/Appendix/danhgiatv.tex}{
    \include{Project1/Appendix/danhgiatv}
  }{}
}

% Phần đánh giá mức độ hoàn thành
\IfFileExists{Appendix/danhgiamucdohoanthanh.tex}{
  \chapter*{Đánh giá mức độ hoàn thành từng danh mục}
\phantomsection
\addcontentsline{toc}{chapter}{Đánh giá mức độ hoàn thành từng danh mục}

Nhóm đã tiến hành rà soát, đối chiếu toàn bộ kết quả thực hiện với các yêu cầu đề ra của Đồ án 1 (Bài toán tự động giải trò chơi Bloxorz bằng các thuật toán tìm kiếm AI). Bảng dưới đây trình bày chi tiết mức độ hoàn thành của từng hạng mục công việc.

\begin{center}
\textbf{Bảng đánh giá chi tiết mức độ hoàn thành theo từng danh mục yêu cầu của đồ án}
\vspace{6pt}

\begingroup
\footnotesize
\setlength{\tabcolsep}{5pt}
\renewcommand{\arraystretch}{1.35}
\begin{tabularx}{\textwidth}{|c|p{0.24\textwidth}|X|c|}
\hline
\rowcolor{gray!15}
\textbf{STT} & \textbf{Danh mục yêu cầu} & \textbf{Nội dung chi tiết và kết quả đạt được} & \textbf{Mức độ hoàn thành} \\
\hline
\multirow{3}{*}{1} & \multirow{3}{=}{\textbf{Mô hình hóa bài toán và xây dựng lõi trò chơi Bloxorz}} 
& Biểu diễn chính xác không gian trạng thái $s = (C, a, b)$ cho cả khối hộp $1 \times 1 \times 2$ (đứng, nằm ngang, nằm dọc) và hai khối đơn $1 \times 1 \times 1$ khi phân tách. & \textbf{100\%} \\
\cline{3-4}
& & Hiện thực đầy đủ toàn bộ cơ chế ô địa hình: ô thường, ô đích ($G$), ô dễ vỡ (\textit{fragile tiles}), công tắc nhẹ hình tròn (\textit{soft switch}), công tắc nặng chữ X (\textit{heavy switch}), công tắc phân tách (\textit{split switch}) và hệ thống cầu đóng/mở theo mặt nạ bit (\textit{bitmask}). & \textbf{100\%} \\
\cline{3-4}
& & Kiểm tra chặt chẽ điều kiện biên, rơi khỏi bàn chơi, sập ô dễ vỡ khi đứng thẳng và xác nhận trạng thái đích hợp lệ. & \textbf{100\%} \\
\hline

\multirow{4}{*}{2} & \multirow{4}{=}{\textbf{Hiện thực các thuật toán tìm kiếm AI}} 
& \textbf{Thuật toán BFS:} Duyệt theo lớp bằng hàng đợi FIFO, kiểm tra trùng lặp trạng thái qua bảng băm, bảo đảm tìm ra lời giải có số bước di chuyển $L$ tối thiểu. & \textbf{100\%} \\
\cline{3-4}
& & \textbf{Thuật toán DFS:} Duyệt theo chiều sâu bằng ngăn xếp LIFO, quản lý chặt chẽ tập trạng thái đã thăm để triệt tiêu hoàn toàn chu trình lặp vô hạn. & \textbf{100\%} \\
\cline{3-4}
& & \textbf{Thuật toán UCS:} Duyệt theo chi phí tích lũy bằng hàng đợi ưu tiên Min-Heap với hàm chi phí địa hình $c(s, u, s') = 1 + 2F + B$, bảo đảm nghiệm có tổng chi phí $C$ nhỏ nhất. & \textbf{100\%} \\
\cline{3-4}
& & \textbf{Thuật toán A*:} Kết hợp chi phí tích lũy $g(n)$ và hàm đánh giá heuristic $h(n)$ dựa trên khoảng cách Chebyshev chia đôi; thỏa mãn tính chấp nhận được (\textit{admissible}) và nhất quán (\textit{consistent}), giúp giảm mạnh số nút mở rộng. & \textbf{100\%} \\
\hline

\multirow{3}{*}{3} & \multirow{3}{=}{\textbf{Giao diện đồ họa (GUI) và trực quan hóa}} 
& Thiết kế giao diện Pygame trực quan với đầy đủ màn hình chính, màn hình chọn màn chơi, hướng dẫn, cài đặt và màn hình chơi trực tiếp. & \textbf{100\%} \\
\cline{3-4}
& & Tích hợp bảng điều khiển AI ngay trong màn hình chơi: chọn thuật toán (BFS, DFS, UCS, A*), diễn hoạt tự động từng bước (\textit{step-by-step replay}), điều chỉnh tốc độ và hiển thị trực tiếp các chỉ số đo lường. & \textbf{100\%} \\
\cline{3-4}
& & Bổ sung tính năng nâng cao: hệ thống đa trang phục (\textit{skins}) cho khối lăn đồng bộ cùng phông nền môi trường, hiệu ứng hạt vật lý (bụi, mảnh vỡ, tia lửa) và âm thanh tổng hợp. & \textbf{100\%} \\
\hline

\multirow{3}{*}{4} & \multirow{3}{=}{\textbf{Thực nghiệm, kiểm thử và đánh giá định lượng}} 
& Thiết kế bộ dữ liệu chuẩn gồm 32 màn kiểm thử (\texttt{*.json}) chia làm 4 nhóm đặc trưng (Cơ bản, Ô dễ vỡ, Công tắc và cầu, Kết hợp) có thống kê cấu trúc chi tiết. & \textbf{100\%} \\
\cline{3-4}
& & Đo đạc tự động 4 chỉ số (Thời gian chạy \texttt{perf\_counter}, Đỉnh bộ nhớ \texttt{tracemalloc}, Số nút mở rộng, Độ dài và chi phí lời giải) qua 128 lượt khởi động và 640 lượt đo chính thức. & \textbf{100\%} \\
\cline{3-4}
& & Xây dựng bộ kiểm thử tự động (\texttt{pytest}) kiểm chứng lại 100\% lời giải bằng phương thức \texttt{Game.step}, bảo đảm tính chính xác tuyệt đối của mọi thuật toán. & \textbf{100\%} \\
\hline

\multirow{3}{*}{5} & \multirow{3}{=}{\textbf{Báo cáo khoa học và đóng gói sản phẩm}} 
& Biên soạn báo cáo chuẩn LaTeX gồm 6 chương và phụ lục: trình bày rõ ràng cơ sở lý thuyết, hình thức hóa toán học, mã giả thuật toán và chứng minh tính chất của hàm heuristic. & \textbf{100\%} \\
\cline{3-4}
& & Trực quan hóa kết quả thực nghiệm bằng hệ thống bảng thống kê đầy đủ và biểu đồ so sánh 4 thuật toán theo từng nhóm màn chơi, rút ra nhận xét chuyên sâu về ưu/nhược điểm. & \textbf{100\%} \\
\cline{3-4}
& & Hoàn thiện danh mục tài liệu và mã nguồn tham khảo, tổ chức kho mã nguồn khoa học kèm tài liệu hướng dẫn cài đặt và video demo. & \textbf{100\%} \\
\hline
\rowcolor{gray!15}
\multicolumn{3}{|r|}{\textbf{TỔNG ĐÁNH GIÁ MỨC ĐỘ HOÀN THÀNH TOÀN ĐỒ ÁN}} & \textbf{100\%} \\
\hline
\end{tabularx}
\endgroup
\end{center}

}{
  \IfFileExists{Project1/Appendix/danhgiamucdohoanthanh.tex}{
    \include{Project1/Appendix/danhgiamucdohoanthanh}
  }{}
}

% Tạo danh sách hình bảng
\listoffigures
\addcontentsline{toc}{chapter}{Danh sách hình}
\listoftables
\addcontentsline{toc}{chapter}{Danh sách bảng}

\IfFileExists{Appendix/noidung.tex}{\include{Appendix/noidung}}{
  \IfFileExists{Project1/Appendix/noidung.tex}{\include{Project1/Appendix/noidung}}{}
}

\clearpage

\pagenumbering{arabic} % Đánh số 1, 2, 3, ...

% Chương 1: Giới thiệu
\chapter{Giới thiệu}
\label{chap:gioi-thieu}

\section{Bài toán Griductive}
\label{sec:bai-toan-griductive}
% Giới thiệu bối cảnh trò chơi Griductive, quy tắc suy luận tất định và nguyên tắc không đoán mò.

\section{Mục tiêu và phạm vi}
\label{sec:muc-tieu-pham-vi}
% Trình bày các mục tiêu chính của đồ án và phạm vi nghiên cứu, hiện thực hệ thống.

 
% Chương 2: Tổng quan cài đặt trò chơi
\chapter{Tổng quan cài đặt trò chơi}
\label{chap:tong-quan-cai-dat}

\section{Luật chơi và các loại gợi ý}
\label{sec:luat-choi-cac-loai-goi-y}
% Trình bày luật chơi cơ bản, cấu trúc lưới, nhân vật và phân loại các mẫu gợi ý.

\section{Định dạng câu đố và kiến trúc mã nguồn}
\label{sec:dinh-dang-kien-truc}

\subsection{Định dạng câu đố}
\label{subsec:dinh-dang-cau-do}
% Mô tả lược đồ dữ liệu và cấu trúc tệp biểu diễn câu đố.

\subsection{Kiến trúc mã nguồn}
\label{subsec:kien-truc-ma-nguon}
% Trình bày cấu trúc thư mục và phân định vai trò giữa lõi trò chơi và tác tử logic.

\section{Giao diện đồ họa}
\label{sec:giao-dien-do-hoa}

\subsection{Bàn cờ và hiển thị gợi ý}
\label{subsec:ban-co-hien-thi}
% Trình bày thiết kế giao diện bàn cờ và trực quan hóa vùng tham chiếu của gợi ý.

\subsection{Tương tác lật thẻ và phán đoán}
\label{subsec:tuong-tac-lat-the}
% Trình bày tương tác người dùng khi đưa ra phán quyết và phản hồi lật thẻ.

\subsection{Bảng điều khiển bộ giải}
\label{subsec:bang-dieu-khien-bo-giai}
% Mô tả các chức năng điều khiển nạp màn, khởi động lại, gợi ý nước đi và tự động giải.

\section{Video demo}
\label{sec:video-demo}
% Giới thiệu video minh họa hoạt động của chương trình và liên kết xem video.

 
% Chương 3: Hình thức hóa bài toán suy luận
\chapter{Hình thức hóa bài toán suy luận}
\label{chap:hinh-thuc-hoa}

\section{Biểu diễn biến mệnh đề và cơ sở tri thức}
\label{sec:bien-menh-de-kb}

Để áp dụng các kỹ thuật suy luận hình thức vào việc giải quyết trò chơi Griductive, toàn bộ trạng thái của bàn cờ và nội dung các gợi ý được chuyển đổi sang ngôn ngữ của logic mệnh đề. 

Xét một bàn cờ kích thước $n \times n$ với $n$ nhận giá trị thuộc tập 3, 4 hoặc 5. Tọa độ của mỗi ô cờ được chuẩn hóa thành cặp chỉ số gồm chỉ số hàng $r$ và chỉ số cột $c$, trong đó cả hai chỉ số đều nhận giá trị nguyên từ 0 đến $n - 1$. Với mỗi vị trí ô cờ $i$ tương ứng với tọa độ $(r, c)$, ta định nghĩa đúng một biến mệnh đề Boole cơ sở ký hiệu là $C_i$. Giá trị chân lý của biến này phản ánh thân phận bí mật của nhân vật tại ô cờ tương ứng:
\[
C_i = 
\begin{cases}
\text{True}, & \text{nếu nhân vật tại ô } i \text{ là Tội phạm}, \\
\text{False}, & \text{nếu nhân vật tại ô } i \text{ là Vô tội}.
\end{cases}
\]

Nhằm bảo đảm tính nhất quán trong toàn bộ quá trình mã hóa và giải toán, ánh xạ từ tọa độ ô cờ sang tập các biến số nguyên dương được thiết lập theo thứ tự ưu tiên hàng trước như sau:
\[
\text{var}(r, c) = r \cdot n + c + 1 \in \{1, 2, \dots, n^2\}
\]

Hàm ánh xạ trên thiết lập một quan hệ song ánh tất định giữa $n^2$ ô cờ và tập các số nguyên từ 1 đến $n^2$. Đối với các mệnh đề đếm số lượng phức tạp đòi hỏi sử dụng thêm các biến phụ nhằm giảm thiểu số lượng mệnh đề tuyển, các biến phụ này luôn được cấp phát bắt đầu từ chỉ số $n^2 + 1$ trở đi và được phân lập rạch ròi khỏi tập các biến mệnh đề cơ sở.

Tại mỗi bước suy luận $t$ trong tiến trình trò chơi, cơ sở tri thức công khai ký hiệu là $KB_t$ đại diện cho toàn bộ thông tin hợp thức mà người chơi hoặc tác tử logic được phép tiếp cận. Cơ sở tri thức này được kiến tạo từ hai nguồn dữ kiện chính: tập hợp các mệnh đề chuẩn hội chuyển đổi từ các gợi ý đã được lật mở và tập hợp các mệnh đề đơn vị khẳng định thân phận của các nhân vật đã được chứng minh:
\[
KB_t = \left( \bigwedge_{c \in \text{RevealedClues}_t} \text{CNF}(c) \right) \land \left( \bigwedge_{i \in \text{KnownStatuses}_t} \text{UnitClause}(i, \text{Status}(i)) \right)
\]

Cơ sở tri thức $KB_t$ phải tuân thủ nghiêm ngặt nguyên tắc bảo toàn thông tin: tuyệt đối không chứa thân phận ẩn của các ô cờ đang úp, không chứa nội dung của các gợi ý chưa được lật, và không chứa bất kỳ giả định bổ sung nào ngoài những dữ kiện đã công bố.

Tiến trình lật thẻ bài được điều phối thông qua giao thức suy luận dựa trên nguyên lý hệ quả logic. Trước tiên, hệ thống kiểm tra tính thỏa được của cơ sở tri thức công khai $KB_t$. Nếu $KB_t$ rơi vào trạng thái không thỏa được, điều đó đồng nghĩa với việc xuất hiện mâu thuẫn trong tập dữ kiện và hệ thống sẽ thông báo trạng thái mâu thuẫn. Khi $KB_t$ thỏa được, với mỗi ô cờ $i$ chưa lật, việc xác định thân phận của nhân vật được thực hiện thông qua phương pháp phản chứng:
\[
KB_t \models C_i \iff KB_t \land \neg C_i \text{ không thỏa được}
\]
\[
KB_t \models \neg C_i \iff KB_t \land C_i \text{ không thỏa được}
\]

Dựa trên kết quả kiểm tra tính thỏa được của hai công thức trên, trạng thái suy luận của ô cờ $i$ được phân loại cụ thể như sau:
\begin{itemize}
    \item Nếu $KB_t \land \neg C_i$ không thỏa được, nhân vật tại $i$ bắt buộc phải là Tội phạm.
    \item Nếu $KB_t \land C_i$ không thỏa được, nhân vật tại $i$ bắt buộc phải là Vô tội.
    \item Nếu cả hai công thức $KB_t \land \neg C_i$ và $KB_t \land C_i$ đều thỏa được, thông tin hiện tại chưa đủ để đưa ra kết luận và trạng thái của ô cờ được xác định là Chưa xác định.
\end{itemize}

Khi người chơi hoặc tác tử gửi một phán quyết về thân phận của ô cờ $i$, hệ thống sẽ phản hồi theo một trong ba trường hợp: Chấp nhận phán quyết khi và chỉ khi thân phận đó đã được chứng minh bắt buộc từ cơ sở tri thức hiện tại, khi đó thẻ bài được lật mở và gợi ý mới được bổ sung vào $KB_{t+1}$; Từ chối với lý do chưa thể chứng minh khi ô cờ mang trạng thái chưa xác định; và Từ chối với lý do mâu thuẫn khi ô cờ bị suy diễn bắt buộc mang thân phận ngược lại với phán quyết. Trong hai trường hợp từ chối, trạng thái bàn cờ hoàn toàn không thay đổi và thẻ bài vẫn giữ nguyên trạng thái úp nhằm ngăn chặn hành vi đoán mò.

\section{Ngữ nghĩa các mẫu gợi ý}
\label{sec:ngu-nghia-goi-y}

\subsection{Các mẫu gợi ý cơ bản}
\label{subsec:goi-y-co-ban}

Hệ thống hỗ trợ sáu mẫu gợi ý cốt lõi, phản ánh các mối quan hệ logic cơ bản giữa các nhân vật và các vùng hình học trên lưới:

\begin{enumerate}
    \item Gợi ý sự thật đơn lẻ: Xác nhận trực tiếp thân phận của đúng một ô cờ $i$. Biểu thức logic tương ứng là mệnh đề đơn vị $C_i$ nếu nhân vật là Tội phạm, hoặc phủ định $\neg C_i$ nếu nhân vật là Vô tội.
    \item Gợi ý đồng nhất trạng thái: Xác nhận hai nhân vật tại ô $i$ và ô $j$ có cùng thân phận với nhau, nghĩa là cả hai cùng là Tội phạm hoặc cùng là Vô tội. Mối quan hệ này tương ứng với phép tương đương logic:
    \[
    C_i \leftrightarrow C_j
    \]
    \item Gợi ý khác biệt trạng thái: Xác nhận hai nhân vật tại ô $i$ và ô $j$ có thân phận đối lập nhau, nghĩa là một người là Tội phạm và người còn lại là Vô tội. Mối quan hệ này tương ứng với phép tuyển loại trừ:
    \[
    C_i \oplus C_j
    \]
    \item Gợi ý đếm chính xác số lượng: Khẳng định trong vùng tham chiếu $R$, tổng số lượng tội phạm đúng bằng một số nguyên $k$, với điều kiện $k$ nằm trong đoạn từ 0 đến kích thước của vùng $R$. Biểu thức toán học tương ứng là:
    \[
    \sum_{i \in R} C_i = k
    \]
    \item Gợi ý đếm số lượng tối thiểu: Khẳng định trong vùng tham chiếu $R$, số lượng tội phạm đạt ít nhất là $k$, tương ứng với bất đẳng thức:
    \[
    \sum_{i \in R} C_i \ge k
    \]
    \item Gợi ý đếm số lượng tối đa: Khẳng định trong vùng tham chiếu $R$, số lượng tội phạm không vượt quá $k$, tương ứng với bất đẳng thức:
    \[
    \sum_{i \in R} C_i \le k
    \]
\end{enumerate}

\subsection{Các mẫu gợi ý mở rộng}
\label{subsec:goi-y-mo-rong}

Bên cạnh sáu mẫu cốt lõi, đồ án hiện thực thêm hai mẫu gợi ý mở rộng nhằm tăng cường khả năng biểu diễn các ràng buộc phức tạp hơn trong thực tế:

\begin{enumerate}
    \item Ràng buộc tính chẵn lẻ: Xác định tính chất chẵn lẻ của tổng số tội phạm có mặt trong một vùng tham chiếu $R$. Biểu thức toán học được mô tả thông qua phép toán đồng dư theo mô-đun 2:
    \[
    \left(\sum_{i \in R} C_i\right) \bmod 2 = 
    \begin{cases}
    0, & \text{nếu tính chất là Chẵn}, \\
    1, & \text{nếu tính chất là Lẻ}.
    \end{cases}
    \]
    \item Ràng buộc so sánh tương đối giữa hai vùng: Thiết lập mối quan hệ định lượng giữa số lượng tội phạm trong vùng bên trái $L$ và vùng bên phải $R$ thông qua một toán tử so sánh thuộc tập các phép toán lớn hơn, lớn hơn hoặc bằng, nhỏ hơn, nhỏ hơn hoặc bằng, bằng nhau hoặc khác nhau:
    \[
    \sum_{i \in L} C_i \quad \text{op} \quad \sum_{j \in R} C_j
    \]
    Trong trường hợp hai vùng tham chiếu có các ô cờ giao nhau, các ô thuộc phần giao này được tính độc lập ở cả hai vế theo đúng định nghĩa số học, bảo đảm tính tự nhiên và chính xác của phát biểu logic mà không yêu cầu người tạo câu đố phải khử trùng lặp thủ công.
\end{enumerate}

\section{Mã hóa sang dạng chuẩn hội (CNF)}
\label{sec:ma-hoa-cnf}

\subsection{Dẫn xuất toán học cho gợi ý nhị phân}
\label{subsec:dan-xuat-nhi-phan}
% Phần này do thành viên phụ trách bộ mã hóa (Module B) biên soạn.

\subsection{Dẫn xuất toán học cho gợi ý đếm số lượng}
\label{subsec:dan-xuat-dem-so-luong}
% Phần này do thành viên phụ trách bộ mã hóa (Module B) biên soạn.

\subsection{Mã hóa các gợi ý mở rộng và quản lý biến phụ}
\label{subsec:ma-hoa-mo-rong-bien-phu}
% Phần này do thành viên phụ trách bộ mã hóa (Module B) biên soạn.

\section{Đánh giá ngữ nghĩa trực tiếp}
\label{sec:danh-gia-ngu-nghia-truc-tiep}
% Phần này do thành viên phụ trách bộ mã hóa (Module B) biên soạn.

 
% Chương 4: Các thuật toán suy luận
\chapter{Các thuật toán suy luận}
\label{chap:cac-thuat-toan}

\section{Thuật toán DPLL giải bài toán SAT}
\label{sec:thuat-toan-dpll}
% Trình bày thủ tục DPLL với lan truyền đơn vị, phát hiện mâu thuẫn, chọn biến và quay lui.

\section{Thuật toán suy diễn logic và phân loại trạng thái}
\label{sec:suy-dien-logic}
% Trình bày nguyên lý hệ quả logic qua bài toán không thỏa được và quy trình phân loại trạng thái.

\section{Thuật toán kiểm tra tính duy nhất của nghiệm}
\label{sec:kiem-tra-tinh-duy-nhat}
% Trình bày thuật toán xác định câu đố có nghiệm duy nhất trên toàn bộ tập gợi ý.

\section{Vòng lặp suy luận của tác tử logic}
\label{sec:vong-lap-suy-luan}
% Trình bày vòng lặp suy luận không đoán mò, cập nhật cơ sở tri thức và cơ chế gợi ý nước đi.

 
% Chương 5: Thực nghiệm và đánh giá
\chapter{Thực nghiệm và đánh giá}
\label{chap:thuc-nghiem-danh-gia}

\section{Thiết kế bộ dữ liệu câu đố}
\label{sec:thiet-ke-du-lieu}
% Mô tả phương pháp xây dựng, phân loại và cấu trúc các bộ câu đố kiểm thử.

\section{Phương pháp và chỉ số đo lường}
\label{sec:phuong-phap-do-luong}
% Trình bày môi trường thực nghiệm, quy trình kiểm thử và các chỉ số đo lường hiệu năng.

\section{Kết quả định lượng}
\label{sec:ket-qua-dinh-luong}

\subsection{Kết quả trên các câu đố 3x3}
\label{subsec:ket-qua-3x3}
% Bảng số liệu và nhận xét kết quả trên nhóm câu đố 3x3.

\subsection{Kết quả trên các câu đố 4x4}
\label{subsec:ket-qua-4x4}
% Bảng số liệu và nhận xét kết quả trên nhóm câu đố 4x4.

\subsection{Kết quả trên câu đố 5x5}
\label{subsec:ket-qua-5x5}
% Bảng số liệu và nhận xét kết quả trên câu đố thử tải 5x5.

\section{Phân tích và nhận xét}
\label{sec:phan-tich-nhan-xet}
% Phân tích hiệu năng bộ giải, số biến, số mệnh đề, các trường hợp bế tắc hoặc mâu thuẫn.

\section{Giới hạn thực nghiệm}
\label{sec:gioi-han-thuc-nghiem}
% Nêu các giới hạn trong thiết lập thực nghiệm và phạm vi đánh giá của đồ án.

 
% Chương 6: Kết luận và hướng phát triển
\chapter{Kết luận và hướng phát triển}
\label{chap:ket-luan-huong-phat-trien}

\section{Tổng kết đề tài}
\label{sec:tong-ket-de-tai}
% Tóm lược toàn bộ quá trình thực hiện đồ án và các thành phần đã xây dựng.

\section{Kết quả đạt được}
\label{sec:ket-qua-dat-duoc}
% Đánh giá các kết quả cụ thể về mô hình logic, thuật toán, thực nghiệm và sản phẩm ứng dụng.

\section{Hạn chế của đề tài}
\label{sec:han-che-de-tai}
% Nêu các khó khăn, thách thức và hạn chế chưa giải quyết trọn vẹn trong đồ án.

\section{Hướng phát triển trong tương lai}
\label{sec:huong-phat-trien}
% Đề xuất các giải pháp cải tiến, tối ưu hóa thuật toán và mở rộng tính năng của hệ thống.


\chapter{Phụ lục}

\section{Prompt cải thiện giao diện}
\label{app:updateui}
 
\subsection*{Prompt: Thêm hiệu ứng, âm thanh}
\label{sec:prompt-u1}

\begin{quote}
\small
\setlength{\emergencystretch}{2em}
Thêm lớp cảm giác cho game Bloxorz (pygame-ce) mà không đổi luật chơi.

\begin{enumerate}
  \item Hiệu ứng hạt, mô phỏng trong tọa độ thế giới (hàng, cột, độ cao) và chiếu theo góc nhìn của bàn:
  \begin{itemize}
    \item Bụi: bắn ra từ rìa vùng chân khối khi khối đáp xuống; vẽ phía sau khối.
    \item Mảnh vỡ: bay lên từ ô dễ vỡ khi ô gãy; vẽ phía trước khối.
    \item Tia lửa: bật ra từ công tắc vừa được kích hoạt, bắt đầu ngay ngoài vùng khối đang đứng trên nó; vẽ phía trước khối và có hiệu ứng cộng sáng.
    \item Số hạt tồn tại cùng lúc phải bị giới hạn ở mức tối đa.
  \end{itemize}

  \item Âm thanh được tổng hợp hoàn toàn bằng mã, không dùng tệp âm thanh bên ngoài, gồm các âm: lăn, đổi khối, công tắc, vỡ, rơi, va chạm, thắng.

  \item Nâng cấp hình ảnh bàn và khối:
  \begin{itemize}
    \item Mặt đá sáng hơn so với cạnh bên, cạnh mảnh có bevel.
    \item Có bóng tiếp xúc dưới khối.
    \item Ba mặt của khối phân cấp độ sáng rõ ràng.
  \end{itemize}

  \item Mỗi sự kiện trong game (đáp, vỡ, bật công tắc, rơi, thắng) phải kích hoạt hiệu ứng và âm thanh tương ứng.
\end{enumerate}
\end{quote}

\subsection*{Prompt: Đổi phong cách giao diện}
\label{sec:prompt-u4}

\begin{quote}
\small
\setlength{\emergencystretch}{2em}
Nâng cấp trang chủ và không khí hình ảnh của Bloxorz theo phong cách hầm ngục cổ (obsidian và vàng rèn).

\begin{enumerate}
  \item Lớp không khí hình ảnh, vẽ theo thứ tự chiều sâu từ xa đến gần:
  \begin{itemize}
    \item Ảnh nền có hiệu ứng parallax theo chuyển động của chuột.
    \item Tia sáng xiên, sương trôi, đốm lửa và bụi lơ lửng, sao băng.
    \item Sóng lan ra khi khối đáp xuống.
    \item Viền tối (vignette) quanh màn hình.
    \item Mưa sao băng ăn mừng khi thắng.
  \end{itemize}

  \item Trang chủ:
  \begin{itemize}
    \item Chữ BLOXORZ gồm 7 khối đá 3D (chữ đùn khối, có bóng), rơi tuần tự theo trọng lực. Mỗi chữ chạm đáy sinh tia lửa, vòng sóng và tiếng động; chữ cuối kèm rung nhẹ.
    \item Bên phải là bàn cờ mẫu có khối $1\times1\times2$ tự lăn, chu kỳ đồng bộ với hiệu ứng của chữ và lặp lại khi chữ khởi động lại.
    \item Nút chính có chevron vàng trượt khi di chuột lên.
    \item Các animation không được chặn việc bấm nút.
  \end{itemize}

  \item Camera 3D:
  \begin{itemize}
    \item Có 5 chế độ góc nhìn: Cổ điển ($25^\circ$/$-20^\circ$), Isometric ($30^\circ$/$+45^\circ$), Góc phải ($28^\circ$/$+32^\circ$), Trên cao ($68^\circ$/$-15^\circ$), Điện ảnh ($15^\circ$/$-35^\circ$).
    \item Chuyển giữa các góc nhìn mượt; phím V đổi góc nhìn.
    \item Phím Q/E hoặc kéo chuột phải để xoay quanh bàn (orbit).
    \item Góc màn hình có la bàn khối hộp 3D thu nhỏ.
  \end{itemize}

  \item Hiệu ứng thị giác bổ sung:
  \begin{itemize}
    \item Xoáy vực và ảnh mờ dư khi khối rơi.
    \item Rune nổi trên khối đang được chọn khi khối tách.
    \item Mảnh rune bay ra khi thắng.
  \end{itemize}
\end{enumerate}
\end{quote}

\subsection*{Prompt: Nâng cấp phong cách giao diện hoàn chỉnh}
\label{sec:prompt-obsidian}

\begin{quote}
\small
\setlength{\emergencystretch}{2em}
Nâng cấp giao diện Bloxorz lên bản Obsidian Sanctuary hoàn chỉnh. Giữ nguyên luật chơi và tiêu chí của các solver BFS/DFS/UCS/A*. Làm lần lượt theo từng cụm.

\begin{enumerate}[label=\Alph*.]
  \item Bố cục thích ứng:
  \begin{itemize}
    \item Canvas có đúng kích thước vùng hiển thị của cửa sổ, không scale từ $1920\times1080$.
    \item Từ 1100~px trở lên: bàn ở bên trái, thanh bên (sidebar) ở bên phải rộng 400~px (có thể 450 hoặc 500~px ở màn rộng).
    \item Dưới 1100~px: thanh bên trở thành ngăn AI cuộn được, mở bằng nút hoặc phím F2, tự đóng khi bắt đầu phát lời giải.
    \item Thay đổi kích thước cửa sổ phải giữ nguyên trạng thái game, lịch sử, solver, camera và kết quả.
    \item Kiểm tra ở các kích thước $768\times432$, $1280\times720$, $1600\times900$, $1920\times1080$, $2880\times1800$.
  \end{itemize}

  \item Floating Dock: thanh kính mờ bo tròn ở giữa đỉnh màn hình, hiển thị tên màn, số bước, trạng thái và nút chính; tự thu gọn thành capsule khi khối đang lăn, bung ra khi rê chuột.

  \item Chiến dịch:
  \begin{itemize}
    \item 33 màn chia thành 4 giai đoạn (1--10, 11--20, 21--30, 31--33), mỗi giai đoạn có tên và mô tả riêng.
    \item Banner điện ảnh mở đầu mỗi màn.
    \item Nút Tiếp tục/Bắt đầu dựa trên màn chiến dịch đã lưu trong tùy chọn người chơi.
    \item Hộp thoại chọn chế độ: phím 1 chọn Chiến dịch, phím 2 chọn Sandbox, Esc để đóng.
    \item Không khóa màn nếu chưa có yêu cầu.
  \end{itemize}

  \item Trang phục (skin) của khối:
  \begin{itemize}
    \item Mỗi skin gồm màu gốc, hệ số khung/viền/sheen, kiểu rivet, glyph trung tâm, màu tia lửa, màu khối tách, bóng gợi ý và texture tùy chọn.
    \item Có 7 skin dựng sẵn: classic, magma, cyber, ice, jade, gold, wood.
    \item Chọn skin đổi cả bảng màu của menu và của bàn.
  \end{itemize}

  \item Cài đặt:
  \begin{itemize}
    \item Hộp thoại (modal) có 4 tab: Âm thanh, Hiển thị, Bàn chơi, Trang phục. Tab Trang phục có xem trước khối 3D xoay.
    \item Tùy chọn người chơi được lưu dưới dạng JSON theo cách nguyên tử (ghi tệp tạm rồi thay thế); đọc lỗi thì dùng giá trị mặc định.
    \item Mỗi tùy chọn (âm thanh, nhạc nền, toàn màn hình, giảm chuyển động, hạt nền, tọa độ ô, đường lời giải) phải tác động thật lên game.
  \end{itemize}

  \item Chọn màn:
  \begin{itemize}
    \item Ô tìm kiếm riêng, hỗ trợ nhập văn bản, Ctrl+F, Ctrl+A, Esc, và tìm không dấu theo tên hoặc số màn.
    \item Lọc theo nguồn màn và theo cơ chế (suy ra từ các loại ô và công tắc có trong màn).
    \item Thanh công cụ cố định; thẻ màn bị cắt ở mép dưới vùng cuộn.
    \item Bố cục 4/3/2 cột tùy độ rộng.
    \item Thumbnail dựng lười, chỉ dựng cho thẻ đang hiển thị.
    \item Giữ lựa chọn và bộ lọc khi đổi kích thước cửa sổ.
  \end{itemize}

  \item Trải nghiệm solver:
  \begin{itemize}
    \item Các trạng thái hiển thị: chờ, đang tìm (kèm thời gian), gợi ý, phát, tạm dừng, thành công, vô nghiệm, hết giới hạn, đã hủy, lỗi.
    \item Tìm kiếm chạy nền; khi đổi màn hoặc hủy thì kết quả cũ bị bỏ.
    \item Giới hạn mặc định: 30~s hoặc 500.000 nút.
    \item Phát lại: tạm dừng/tiếp tục, bước tới, tốc độ $0.5$--$3\times$.
    \item Số bước của AI được tách khỏi tổng số bước của lượt chơi.
  \end{itemize}

  \item So sánh solver:
  \begin{itemize}
    \item Bốn solver chạy tuần tự trên cùng một trạng thái game được chụp tại thời điểm người chơi bấm.
    \item Có hai tab Kết quả và Biểu đồ.
    \item Sắp xếp theo số bước, chi phí, thời gian, số nút, bộ nhớ; thứ tự mặc định là BFS, DFS, UCS, A*.
    \item Chỉ tô giá trị nhỏ nhất trong các kết quả thành công.
  \end{itemize}
\end{enumerate}
\end{quote}

\subsection*{Prompt: Nền riêng cho từng skin}
\label{sec:prompt-u8}

\begin{quote}
\small
\setlength{\emergencystretch}{2em}
Thêm mỗi skin có ảnh nền riêng.

\begin{enumerate}
  \item Mỗi skin có thêm các thuộc tính: ảnh nền, màu pha khí quyển, màu tia sáng xiên và màu viền tối (vignette). Giá trị mặc định là bộ nền hiện có.

  \item Tải và lưu đệm (cache) ảnh nền theo từng skin; nếu thiếu tệp ảnh của skin thì dùng nền mặc định. Việc dựng nền chỉ thực hiện một lần cho mỗi tổ hợp kích thước, skin và độ phủ tối, và được dựng lại khi các thông số này thay đổi. Bộ đệm bị xóa khi đổi skin hoặc đổi kích thước cửa sổ.

  \item Độ phủ tối mặc định là 36 (không phải 92), để tranh nền hiện rõ nhưng bàn cờ vẫn là tiêu điểm.

  \item Lớp không khí hình ảnh lấy màu tia sáng và màu viền tối từ skin đang chọn.

  \item Tạo 6 ảnh nền JPG kích thước $1920\times1080$, tối ưu dung lượng, cho các skin cyber, gold, ice, jade, magma và wood. Kèm theo công cụ sinh và cài đặt các ảnh này, và công cụ kiểm tra sự tồn tại, kích thước và độ sáng trung bình của từng ảnh.
\end{enumerate}
\end{quote}

\section{Prompt thiết kế bộ màn kiểm thử}
\label{app:testcase}
 
Prompt dưới đây được sử dụng để xây dựng bộ 32 màn kiểm thử Bloxorz.
 
\begin{quote}
\small
\setlength{\emergencystretch}{2em}
Bạn là người thiết kế dữ liệu kiểm thử cho game Bloxorz trong repository này. Hãy đọc bộ nạp JSON, luật lăn khối, ô dễ vỡ và công tắc của game trước khi tạo dữ liệu. Tạo 32 màn mới, chia thành đúng 4 nhóm, mỗi nhóm đúng 8 màn:
 
\noindent A. simple --- Simple cases without advanced tiles. Chỉ dùng sàn thường \#, vị trí xuất phát S, đích G và vực. Không có ô dễ vỡ F, công tắc mềm o, công tắc nặng x, công tắc tấch p hoặc cầu 0-9. Đa dạng bàn rộng, đường hẹp, lỗ khuyết, ngã rẽ, đường vòng và vị trí cần đổi hướng khối. Mục tiêu: kiểm thử chuyển tư thế đứng/nằm, biên bàn, chu trình và tìm đích.
 
\noindent B. fragile --- Cases with fragile tiles. Dùng các ô của nhóm A ở trên và thêm ô dễ vỡ F; không có công tắc, cầu hoặc tách khối. Khối được nằm trên F nhưng đứng trên F sẽ thua; F không phải sàn sụp sau mỗi lượt đi. Bắt buộc phải chiếm ít nhất một F trên mọi lời giải. Đa dạng dải F, mật độ F, chỗ quay khối và đường vòng để phân biệt đường ngắn nhất với đường có chi phí địa hình thấp.
 
\noindent C. switches --- Cases with switches. Dùng \#, S, G, ., công tắc mềm o, công tắc nặng x và cầu nhóm 0--9. Không có F hoặc tách khối p. Công tắc mềm kích hoạt khi nằm hoặc đứng; công tắc nặng chỉ kích hoạt khi đứng. Phối hợp toggle (đảo đóng/mở) và permanent (mở, kích hoạt lại vẫn mở), từ một tới nhiều nhóm cầu. Cầu ban đầu đóng. Mỗi công tắc điều khiển một nhóm cầu cần thiết; vô hiệu hóa riêng bất kỳ công tắc nào phải làm màn vô nghiệm. Mục tiêu: kiểm thử điều kiện kích hoạt và trạng thái gồm vị trí, tư thế, bitmask cầu; cùng vị trí nhưng bitmask khác không được gộp thành một state.
 
\noindent D. mixed --- Cases combining both features. Có cả F, công tắc mềm/nặng và cầu; không có tách khối p. Phải đi qua F và kích hoạt các công tắc cần thiết để tới G. Đa dạng thứ tự mở cầu, vị trí dải F, chỗ cần đứng kích hoạt x và chỗ chỉ được nằm trên F. Mục tiêu: kiểm thử tương tác hai cơ chế.
 
\noindent Yêu cầu chung:
\begin{enumerate}
  \item Mỗi nhóm có 2 màn dễ, 2 trung bình, 2 khó, 2 rất khó. Dùng độ dài lời giải BFS ngắn nhất L: dễ 4--8 bước; trung bình 9--14; khó 15--22; rất khó 23--40.
 
  \item Các màn có hình dạng, kích thước, mật độ lỗ và vị trí S/G khác nhau; không dùng bản sao chỉ đổi tên, dịch tọa độ, xoay hoặc phản chiếu. Tất cả có đúng một S, một G và lời giải hợp lệ theo Game.
 
  \item Lưu 32 JSON trong src/levels/testcase/, category="testcase", tên Testcase Simple/\allowbreak{}Fragile/\allowbreak{}Switches/\allowbreak{}Mixed 01--08, ID riêng không trùng. Giữ lời giải, thống kê và metadata báo cáo ở ngoài thư mục màn chơi.
 
  \item JSON dùng schema sẵn có: name, id, category, description, grid, bridges.\allowbreak{}initially\_\allowbreak{}open và switches[\{pos, kind, targets, mode\}]. Tọa độ [row, col] bắt đầu từ 0. Các hàng grid phải có cùng chiều dài.
 
  \item Tích hợp một tab riêng tên chính xác "testcase" trong thư viện màn chơi, cho phép xem đủ 32 màn, tìm theo nhóm, mở chơi và chạy các solver hiện có.
 
  \item Kiểm chứng bằng BFS thật của game và phát lại từng lời giải tới WIN. 
  % Nhóm B/D: tìm kiếm với ràng buộc cấm chiếm mọi F phải vô nghiệm. Nhóm C/D: thay riêng từng công tắc bằng sàn thường, xóa định nghĩa tương ứng, giữ cầu ban đầu đóng; BFS phải vô nghiệm. Lưu chứng cứ cùng hash SHA-256 của JSON để có thể kiểm tra lại.
 
  % \item Xuất báo cáo Markdown với mô tả rõ đặc điểm từng nhóm, ví dụ từng màn, và 4 bảng thống kê, mỗi bảng 8 hàng, gồm: mã màn, độ khó, kích thước H\(\times\)W, số ô có mặt sàn, số vực, số sàn thường (kể S), số đích G, số F, số o, số x, số ô cầu, số nhóm cầu, tổng ô nâng cao (F+o+x+cầu), L tối ưu và số nút BFS mở rộng. Tổng ô có mặt sàn bao gồm cầu dù ban đầu đóng; tổng đó cộng số vực phải bằng H\(\times\)W. Xuất thêm CSV cùng số liệu và JSON lời giải/chứng cứ kiểm chứng.
 
  % \item Dùng quy trình có seed cố định, ghi rõ seed/cách tái tạo. Không bịa số liệu hoặc dùng số bước của một lời giải bất kỳ làm số bước tối ưu.
 
\end{enumerate}

\section{Video Demo}
\href{https://www.youtube.com/watch?v=KgBhIWJ7A10}{Đường dẫn đến video demo}

\end{quote}




% In tài liệu tham khảo
\IfFileExists{References/references.bib}{
  \nocite{*}
  \printbibheading[title={Tài liệu tham khảo}]
  \addcontentsline{toc}{chapter}{Tài liệu tham khảo}
  \printbibliography[heading=subbibliography, title={Tài liệu tiếng Việt}, keyword=Viet, resetnumbers=true]
  \DeclareNameAlias{sortname}{last-first}
  \DeclareNameAlias{default}{last-first}
  \printbibliography[heading=subbibliography, title={Tài liệu tiếng Anh}, keyword=English]
  \printbibliography[heading=subbibliography, title={Mã nguồn và Liên kết tham khảo}, keyword=Source]
}{}

% % Báo cáo làm việc
% \IfFileExists{Appendix/baocaolamviec.tex}{
%   \chapter*{Báo cáo làm việc hàng tuần}
\addcontentsline{toc}{chapter}{Báo cáo làm việc hàng tuần}
% }{
%   \IfFileExists{Project1/Appendix/baocaolamviec.tex}{
%     \include{Project1/Appendix/baocaolamviec}
%   }{}
% }

\end{document}
